\vspace{-0.1cm}
\section{Experimental Evaluation}
\label{sec:experiments}
\vspace{-0.2cm}
The goal of our experiments is to demonstrate the efficacy and justify the design of \methodname{} in training LLMs how to self-correct by only training on their own data. To this end, we perform a comparative evaluation of \methodname{} against prior methods that also use self-generated data to train for self-correction, and run several ablation studies on two representative reasoning tasks where error correction is crucial.

\textbf{Tasks.} We mainly focus on reasoning problems in math and coding: \textbf{(a)} math problem solving on MATH~\citep{hendrycksmath2021}, and \textbf{(b)} code generation on MBPP~\citep{austin2021program} and HumanEval~\citep{chen2021evaluating}. We use the following train-test splits in our experiments: \textbf{(1)} \textbf{MATH}: following \cite{lightman2023let}, we augment the MATH training set with 4500 problems from the test set, and report results on the remaining 500 problems (MATH500); and \textbf{(2)} \textbf{Code generation}: we train on MBPP and report results on HumanEval, which does not expose test cases to the model. For all tasks, we use binary rewards during training, indicating whether the model's answer matches the ground truth one (for MATH) or passes all test cases (for coding).

\textbf{Evaluation protocol and metrics.} We report the self-correction accuracy on a number of tasks with two sequential attempts at the problem, i.e., one round of self-correction. For code generation, following the evaluation protocol of \citet{ni2024next}, we also report results on MBPP-R, an offline repair task that requires correcting incorrect first-attempt programs generated from PaLM 2.

\textbf{Models.} For all of our experiments on coding problems, we fine-tune Gemini 1.0 Pro and for MATH, we fine-tune Gemini 1.5 Flash. For all evaluations, we use greedy decoding (i.e. temperature 0), except for inference-compute scaling in Section~\ref{sec:inference} where we set temperature to be 0.7. For all training methods, we attempted to use a fixed budget of model samples and gradient updates, and do not vary hyperparameters such as learning rate and batch size between runs. For all RL runs, we selected checkpoints with the highest training reward, although a small held-out validation set of problems can also be used. 

\begin{table}[]
\centering
\renewcommand{\arraystretch}{1.1}
\caption{\footnotesize{\textbf{Performance of \methodname{} on MATH.} \methodname{} not only attains a higher accuracy at both attempts, but also  provides the most positive self-correction performance $\Delta$\textbf{(t1, t2)}.}}
\vspace{-0.1cm}
\resizebox{.95\textwidth}{!}{\begin{tabular}{l||ccccc}
\toprule
\textbf{Approach} & \textbf{Acc.@t1} & \textbf{Acc.@t2} & \textbf{$\Delta$(t1, t2)} & \textbf{$\Delta^{\mathrm{i} \rightarrow \mathrm{c}}$(t1, t2)} & \textbf{$\Delta^{\mathrm{c} \rightarrow \mathrm{i}}$(t1, t2)} \\ \midrule
Base model & 52.6\%          & 41.4\%          & -11.2\%               & 4.6\%                & 15.8\%               \\
Self-Refine~\citep{madaan2023self}      & 52.8\%          & 51.8\%          & -1.0\%                 & 3.2\%                  & 4.2\%                \\
STaR w/ $\mathcal{D}_{\mathrm{StaR}}^+$~\citep{zelikman2022star}     & 53.6\%          & 54.0\%          & 0.4\%                 & 2.6\%                  & 2.2\%                \\
\pairsft{} w/ $\mathcal{D}_{\mathrm{SFT}}$~\citep{welleck2023generating}       & 52.4\%          & 54.2\%          & 1.8\%                & 5.4\%                & 3.6\%                  \\ \midrule
\methodname{} (Ours)                & \textbf{60.0\%}          & \textbf{64.4\%}          & \textbf{4.4\%} & \textbf{5.8\%} & \textbf{1.4\%} \\ \bottomrule
\end{tabular}}
\vspace{-0.15cm}
\label{tab:math}
\end{table}


\textbf{Evaluation prompts.} We use zero-shot CoT prompting for evaluation on MATH, zero-shot prompting for evaluation on HumanEval, and the canonical three-shot prompt for first-attempt training samples on MBPP~\citep{austin2021program}.
At the second attempt, we utilize an instruction that does not reveal the correctness of the previous answer, but asks the model to attempt to deduce whether a mistake exists in its first attempt response, and if so, potentially rewrite its response. Our full prompts and self-correction instructions can be found in Appendix~\ref{app:prompts}.

\textbf{Baselines \& comparisons.} We compare \methodname{} to relevant prior approaches based on prompting or those that fine-tune only a single model for both solving the task and for revising responses, and only use self-generated data. Specifically, we compare to \textbf{Self-Refine}~\citep{madaan2023self}, a representative prompting-based approach to elicit self-correction behaviors from a model, akin to Reflexion~\citep{shinn2023reflexion}.
Among the fine-tuning based approaches, we compare to \textbf{\pairsft{}} based on the approach from \citet{welleck2023generating}, and multi-turn \textbf{STaR}~\citep{zelikman2022star, singh2023beyond} that fine-tune the model by maximizing log-likelihood respectively on synthetically-paired repair traces (\pairsft{}) and successful repair traces (STaR).


\vspace{-0.2cm}
\subsection{Benchmark Results}
\vspace{-0.2cm}

\textbf{MATH.} Our results are shown in Table~\ref{tab:math}, as well as in Figure~\ref{fig:main}. \methodname{} %
exhibits substantially stronger performance on both direct and self-correction accuracies relative to baselines. Notably, the intrinsic self-correction gain $\Delta$\textbf{(t1, t2)} of 4.4\% is the first significantly positive delta, despite having fewer incorrect problems to correct by virtue of its higher Accuracy@t1. Relative to the base model Gemini 1.5 Flash, \methodname{} improves $\Delta$\textbf{(t1, t2)} by 15.6\%, and Accuracy@t2 by 23.0\%, and over the next best prior approach, \pairsft{}, by 10.2\% and 2.6\% respectively. By observing the frequency of problems that change from incorrect in the first attempt to correct in the second attempt and vice versa, we see that \methodname{} both improves the rate at which it fixes incorrect answers (14.5\%, compared to 9.5\% for base) and reduces the proportion of correct answers it changes (15.8\% to 1.4\%). 

\begin{table}[t]
\centering
\renewcommand{\arraystretch}{1.1}
\caption{\footnotesize{\textbf{Performance of \methodname{} on HumanEval.} \methodname{} attains the highest self-correction performance (\textbf{Accuracy@t2}, \textbf{$\Delta$(t1, t2)}), and also outperforms other methods at offline correction (\textbf{MBPP-R}).}}
\vspace{-0.1cm}
\resizebox{.9\textwidth}{!}{\begin{tabular}{l||cccccc}
\toprule
\textbf{Method}       & \textbf{MBPP-R} & \textbf{Acc.@t1} & \textbf{Acc.@t2} & \textbf{$\Delta$(t1, t2)} & \textbf{$\Delta^{\mathrm{i} \rightarrow \mathrm{c}}$(t1, t2)} & \textbf{$\Delta^{\mathrm{c} \rightarrow \mathrm{i}}$(t1, t2)} \\ \midrule
Base model      & 47.3\%         & 53.7\%      & 56.7\%      & 3.0\%         & 7.9\%         & 4.9\%         \\
Self-Refine     & 30.7\%             & 53.7\%         & 52.5\%         & -1.2\%            & 9.8\%           & 11.0\%           \\
\pairsft{}             & 59.8\%             & 56.1\%         & 54.3\%         & -1.8\%            & 4.3\%           & 6.1\%           \\ \midrule
\methodname{} (Ours)   & \textbf{60.6\%}   & 52.4\%         & \textbf{64.6\%}      & \textbf{12.2\%}              & \textbf{15.2\%}        & \textbf{3.0\%}         \\ \bottomrule
\end{tabular}}
\vspace{-0.1cm}
\label{tab:mbpp}
\end{table}

\textbf{Code generation.} Our results for the code generation task are shown in Table~\ref{tab:mbpp}. Generally, we find that \methodname{} achieves both improved self-correction and offline repair performance. For MBPP-R~\citep{ni2024next}, we find that \methodname{} improves the base model from 47.3\% to 60.6\%, which is comparable to the gap between GPT-3.5 (43\%) and GPT-4 (63.2\%). Despite only training on MBPP, we find that \methodname{} is especially effective at generalizing to HumanEval, achieving a \textbf{12.2\%} intrinsic self-correction delta, or 9\% higher than the base model. By contrast, \pairsft{} works nearly as well on the static repair task MBPP-R, but actually degrades the base model when evaluated in the self-correction setting, thus underscoring the importances of on-policy sampling for self-correction.

\vspace{-0.2cm}
\subsection{Inference-Compute Scaling with Self-Correction}
\label{sec:inference}
\vspace{-0.2cm}
Next, we investigate if \methodname{} can be used in conjunction with inference-time compute scaling strategies. To do so, we evaluate self-consistency decoding~\citep{wang2022self}, where we sample a diverse set of solutions, and then select the most consistent answer among these solutions. Typically, the default strategy is to sample all solutions in parallel to perform majority voting. However, we show in Figure~\ref{fig:main}~(right) that instead of sampling $2K$ solutions in parallel, it is \emph{more} compute-efficient to sample $K$ solutions in parallel, then perform one round of self-correction on each solution. With 32 solution budget per problem, parallel sampling shows a 7.4\% accuracy gain, while combining it with sequential sampling using self-correction yields a 10.5\% improvement.

\vspace{-0.2cm}
\subsection{Ablation Studies: Understanding the Impact of \methodname{} Components}
\vspace{-0.1cm}

\begin{table}[t]
\centering
\renewcommand{\arraystretch}{1.1}
\caption{\footnotesize{\textbf{Ablation studies to understand the impact of various components in \methodname{}.} Observe that while single-turn training is effective at optimizing the first-attempt accuracy of the model, it leads to degradation in the second attempt. The performance improvements without Stage I or without reward shaping in \methodname{} are small when measured by the difference in accuracy over the two attempts. Utilizing STaR generally leads to worse performance even when it is run from an effective Stage I checkpoint.}}
\vspace{-0.2cm}
\resizebox{.9\textwidth}{!}{\begin{tabular}{lccc}
\toprule
\textbf{Method}              & \textbf{Accuracy@t1} & \textbf{Accuracy@t2} & \textbf{$\Delta$(t1, t2)} \\ \midrule
\methodname{} (Ours)               & 60.0\%          & \textbf{64.4\%}          & \textbf{4.4\%} \\ \hdashline
w/o multi-turn training       & \textbf{61.8}\%          & 59.4\%          & -2.4\%         \\ 
w/o Stage I       & 59.2\%          & 61.4\%          & 2.2\%          \\ 
w/o reward shaping            & 60.0\%          & 62.6\%          & 2.6\% \\
w/ STaR instead of REINFORCE Stage II & 56.2\% & 58.4\% & 2.2\%
\\ \bottomrule
\end{tabular}}
\label{tab:math_ablations}
\vspace{-0.1cm}
\end{table}
Finally, we also present a number of ablation studies to understand the importance of various components in \methodname{}. We perform these ablations on the MATH dataset. Concretely, we aim to answer the following questions: \textbf{(1)} \textbf{the importance of multi-turn training:} Can RL trained to maximize single-turn performance achieve better accuracy@t1 or accuracy@t2?; \textbf{(2)} \textbf{the importance of multi-stage training:} How essential is Stage I to \methodname{}? In other words, why not run Stage II directly?; \textbf{(3)} \textbf{the impact of reward shaping.} How would removing the reward shaping terms affect performance of \methodname{} in Stage II, assuming Stage I was done identically?; \textbf{(4)} \textbf{the importance of on-policy RL:} What if we replaced REINFORCE in Stage II with STaR?.

The results of all of these ablation experiments are shown in  Table~\ref{tab:math_ablations}. As expected, single-turn training improves turn 1 performance, but has negative \textbf{$\Delta$(t1, t2)}. As shown in Figure~\ref{fig:naive_multiturn}, Stage I is critical to \methodname{}; without it, the model achieves 2\% lower $\textbf{$\Delta$(t1, t2)}$ and 3\% lower accuracy@t2. Similarly, we find that removing reward shaping also hurts performance, indicating that the RL objectives in both stages play a significant role in teaching the self-correction behavior. We also find that replacing REINFORCE with STaR in Stage II results in significantly lower absolute performance with no visible improvements in self-improvement performance, which contrasts with the findings in \cite{havrilla2024teaching} that STaR and on-policy RL have similar convergence rates for single-turn RL. This suggests that leveraging on-policy samples is especially critical in the self-correction setting, which presents a multi-turn problem that admits potentially spurious solutions.







